% legal.tex
% mainfile: perfbook.tex
% SPDX-License-Identifier: CC-BY-SA-3.0

\section*{法律声明}

本作品代表编辑和作者的观点，不一定代表其各自雇主的观点。\\
\\
商标：
\begin{itemize}
\item	IBM、z~Systems 和 PowerPC 是 International Business Machines Corporation
	在美国和/或其他国家的商标或注册商标。
\item	Linux 是 Linus Torvalds 的注册商标。
\item	Intel、Itanium、Intel Core 和 Intel Xeon 是
	Intel Corporation 或其子公司在美国和/或其他国家的商标。
\item	Arm 是 Arm Limited（或其子公司）在美国和/或其他地区的注册商标。
\item	SPARC 是 SPARC International, Inc. 的注册商标。
	带有 SPARC 商标的产品基于 Sun Microsystems, Inc. 开发的架构。
\item	其他公司、产品和服务名称可能是相应公司的商标或服务标记。
\end{itemize}

本文档中的非源代码文本和图像根据 Creative Commons
Attribution-Share Alike 3.0 United States 许可证的条款提供。\footnote{
	\url{https://creativecommons.org/licenses/by-sa/3.0/us/}}
简而言之，你可以将本文档的内容用于任何目的——个人、商业或其他用途，
只要保留对作者的署名即可。
同样，可以修改本文档、制作衍生作品和翻译版本，
只要这些修改和衍生作品以与原始文档中的非源代码文本和图像相同的条款
向公众提供。

源代码受各种版本的 GPL 保护。\footnote{
	\url{https://www.gnu.org/licenses/gpl-2.0.html}}
其中一些代码仅限 GPLv2，因为它源自 Linux 内核，
而其他代码则为 GPLv2 或更高版本。
请查看 git 仓库\footnote{
	\url{git://git.kernel.org/pub/scm/linux/kernel/git/paulmck/perfbook.git}}
中 CodeSamples 目录下各源文件的注释头以获取确切的许可证信息。
如果你不确定某个代码片段的许可证，应假定为仅限 GPLv2。

合集作品版权 {\textcopyright}~2005--\commityear\ Paul E. McKenney 所有。
每项个人贡献的版权在贡献时归其贡献者所有，
记录于 git 仓库中。
